% Avsnitt 11.8, ur lectures/L11/appendix/b_exercises.md.

\section{Övningar}
\label{c11:sec:exercises}

\begin{aside}
De här övningarna befäster \secref{c11:sec:arch} till \ref{c11:sec:regmap}: blockarkitekturen,
gränssnittet den landar i, och registerkartan den avbildas på. Paketet \code{can_def} som de alla
vilar på skrevs i \chapref{c10:ch}.

Övning~\ref{c11:ex:reason} och \ref{c11:ex:groups} görs på papper. Övning~\ref{c11:ex:entity} och
\ref{c11:ex:iface} görs vid tangentbordet, mot den \file{controller/can_controller.vhd} som skrevs
under det här passet. \code{meta_prev} och dess övningar hör till \chapref{c12:ch}, tillsammans med
\code{bit_timer}.
\end{aside}

% ------------------------------------------------------------------------------------------
\exerciseset{Resonemang om arkitekturen}

\exercise{Vem äger vad}{Resonemang}
\label{c11:ex:reason}
\xpart{a} \code{tx_shift_reg} och \code{rx_shift_reg} vet inte om bitarna de skiftar hör till ID:t,
kontrollfältet, en databyte eller CRC-värdet. Förklara, med hänvisning till blockschemat,
\emph{vilket} block som vet det, och hur det kan återanvända samma skiftregister till alltihop i
stället för att behöva fem olika.

\xpart{b} \code{crc15} matas en bit i taget och får veta när den ska sluta av den av
\code{tx_shift_reg} och \code{rx_shift_reg} som är aktiv för tillfället (via respektive utgångar
\code{stuff} och \code{real_bit_valid}, som grindar \code{crc15.enable}). Varför måste \code{crc15}
få veta att den ska \emph{hoppa över} just stoppbitar, i stället för att bara matas med varje bit som
dyker upp på ledningen?

\xpart{c} \code{can_controller} instansierar de fyra andra blocken inuti sig, så utifrån sett är hela
kontrollern en enda komponent. Anta att någon föreslår att kontrollen av arbitreringsförlust bryts ut
till ett eget block \emph{bredvid} \code{can_controller}. Förklara vad det blocket skulle behöva
tillgång till, och varför det gör lösningen sämre än att låta kontrollen ligga kvar där den är.

\exercise{Att dela upp en annan ram i grupper}{Resonemang}
\label{c11:ex:groups}
\Secref{c11:sec:trace} följer en ram med DLC \code{2} genom blocken och kommer fram till att den når
bussen som åtta laddningar av \code{tx_shift_reg}, tillsammans 50 bitar, plus en 10-bitars svans som
drivs direkt.

\xpart{a} Gör om den uppdelningen för en ram med ID \code{0x7A0} och DLC \code{5}. Ange listan av
grupper (varje grupps innehåll och antal bitar, i sändningsordning), det totala antalet bitar i det
stoppade området, och det totala antalet bitar i ramen innan någon stoppbit satts in.

\xpart{b} Hur många grupper skulle en ram med DLC \code{0} behöva, och vilka poster ur din lista i
\textbf{a)} försvinner? Vad säger det dig om hur \code{can_controller} måste vara strukturerad kring
datafältet, jämfört med identifieraren eller CRC:n?

\xpart{c} Fyra av grupperna i din lista är smalare än 8 bitar, och \code{tx_shift_reg}:s dataingång
är alltid en hel byte. Namnge dem, och säg vad \code{can_controller} måste skicka med tillsammans med
datan för att registret ska sända rätt antal bitar. (Du behöver inte \chapref{c14:ch}:s svar på
\emph{hur}; frågan är vilken information registret inte kan räkna ut på egen hand.)

\xpart{d} Inget av dina gruppantal är antalet bitar som faktiskt drivs ut på ledningen. Förklara vad
som saknas, varför det inte går att veta enbart utifrån DLC:n, och vilket block som ansvarar för det.

% ------------------------------------------------------------------------------------------
\exerciseset{Entiteten och gränssnittet}

\exercise{Entiteten, kontrollerad}{Simulering}
\label{c11:ex:entity}
\xpart{a} Kör \code{make build-project} från roten av gruppens repo, precis som i \chapref{c10:ch}.
Två egna filer bör nu rapporteras, vid sidan av de två utdelade \file{bridge/}-filerna, och varje
testbänk bör fortfarande hoppas över:

\begin{console}
--> controller/can_def.vhd analyzes cleanly
--> controller/can_controller.vhd analyzes cleanly
--> bridge/spi_def.vhd analyzes cleanly
--> bridge/spi_slave.vhd analyzes cleanly

==> meta_prev_tb (skipped: no meta_prev.vhd yet)
==> bit_timer_tb (skipped: no bit_timer.vhd yet)
==> crc15_tb (skipped: no crc15.vhd yet)
==> tx_shift_reg_tb (skipped: no tx_shift_reg.vhd yet)
==> rx_shift_reg_tb (skipped: no rx_shift_reg.vhd yet)
==> can_controller_tb (skipped: no meta_prev.vhd bit_timer.vhd crc15.vhd tx_shift_reg.vhd rx_shift_reg.vhd yet)
==> register_bank_tb (skipped: no register_bank.vhd yet)
==> spi_reg_bridge_tb (skipped: no spi_reg_bridge.vhd yet)

0 testbench(es) run, 8 skipped.
\end{console}

Den andra raden är hela poängen med den här övningen: en entitet med en tom arkitektur är ändå en
komplett, giltig designenhet. Gå nu längre än vad \code{make build-project} gör och elaborera den för
hand, vilket är det andra av GHDL:s tre steg:

\begin{shell}
cd controller
ghdl -a --std=93 can_def.vhd can_controller.vhd
ghdl -e --std=93 can_controller
\end{shell}

Den elaboreras rent, utan några underblock inuti sig och utan något att simulera. Det finns inget
tredje steg att köra här, eftersom en design utan beteende inte har något att kontrollera och ingen
egen testbänk; \code{can_controller_tb} driver den \emph{färdiga} kontrollern och väntar på
\chapref{c17:ch}. Det är i \chapref{c12:ch} det tredje steget dyker upp, med den första modulen som
har en egen testbänk.

\xpart{b} \code{can_controller_tb} ligger redan i \file{controller/}, och den behöver sju moduler,
varav du har skrivit två. Säg vilka fem som saknas och i vilket kapitel var och en kommer. Varför är
det bättre att bygget hoppar över den testbänken än att det försöker och misslyckas?

\xpart{c} \emph{Valfritt, och bara om du har Quartus Prime Lite installerat.} Ta samma fil genom
Quartus, enligt bilaga~\ref{app:quartus}: New Project Wizard, enhet \code{5CEBA4F23C7N}, lägg till
\file{can_def.vhd} och \file{can_controller.vhd}, sätt \code{can_controller} som toppnivåentitet, och
kompilera.

Det går igenom. Läs varningarna i stället för att hoppa över dem, och säg vilka av dem som är
oundvikliga för en arkitektur som inte driver någonting alls. Leta sedan upp antalet pinnar i
fitterns rapport och jämför det med vad ett DE0-CV faktiskt erbjuder en människa: tio omkopplare,
fyra knappar, tio lysdioder. Skriv ner vad du tror måste sitta mellan den här entiteten och kortet,
och spara anteckningen; \chapref{c18:ch} och \ref{c19:ch} bygger precis det, och kallar det
\code{register_bank} och \code{can_spi_node}.

Ingenting senare i boken beror på det här, och \chapref{c12:ch} till \ref{c19:ch} behöver bara GHDL.
Det står här för att just det antalet pinnar är hela argumentet för de två lagren ovanför
kontrollern, sju kapitel i förväg, och att läsa av det ur en riktig fitterrapport gör det till ditt
eget i stället för en siffra som boken påstår.

\xpart{d} Byt tillfälligt plats på \code{tx_bus} och \code{bus_en} i entitetens portlista, så att de
står i omvänd ordning, och kör om \code{make build-project}. Båda är utgångar av typen
\code{std_logic}, så filen analyseras fortfarande rent. Förklara varför ingenting fångar det nu, i
vilket kapitel det först skulle ställa till skada, och hur den skadan skulle se ut på bussen
(ledtråden är \secref{c11:sec:interface}:s avsnitt om bussidans fem portar). Säg sedan varför en
förväxling mellan två portar av samma typ är svårare att hitta än ett felstavat namn skulle ha varit,
givet att allt i den här boken binds positionsvis. Byt tillbaka dem.

\exercise{Att läsa gränssnittet}{Resonemang}
\label{c11:ex:iface}
\xpart{a} En anropare vill sända en ram med ID \code{0x123}, DLC \code{3} och databytesen
\code{AA BB CC}. Lista varje port den måste driva, och värdet på var och en. \code{tx_data} är
\code{data_t}, 64 bitar bred, och byte 0 ligger i de mest signifikanta bitarna: ange den i sin
helhet.

\xpart{b} Samma anropare vill nu veta om ramen blev sänd. Den har tre statusportar till sitt
förfogande (\code{tx_done}, \code{rx_valid}, \code{error}). Vilken pollar den, och vad är den enda
sak den \emph{aldrig} kan få veta av dem, hur ofta den än pollar? (Ledtråden är
\secref{c11:sec:interface}:s anmärkning om vad \code{error} betyder; skälet är
\secref{c10:sec:frameformat}:s ACK-fält.)

\xpart{c} \code{error} ligger hög tills nästa ram börjar, medan \code{tx_done} och \code{rx_valid}
är pulser på en cykel. Förklara vad som skulle gå fel för parallellklassens C++-drivrutin om
\code{error} vore en puls som de andra två.
